\section{Conclusiones y Líneas de Trabajo Futuro}

El objetivo principal de este trabajo ha sido estudiar la aplicación de arquitecturas Transformer al problema de la predicción de descartes en Riichi Mahjong mediante aprendizaje por imitación utilizando partidas reales de Tenhou. Para ello se diseñó una representación propia del estado de juego, un proceso de tokenización adaptado al dominio y un modelo basado en un Transformer Encoder capaz de aprender a seleccionar descartes a partir de ejemplos de jugadores expertos.

Los resultados obtenidos muestran que el modelo es capaz de aprender patrones relevantes para la selección de descartes a partir de los datos disponibles. Sobre un conjunto reservado de 52.563 estados, el modelo entrenado con 100.000 muestras alcanzó una precisión exacta del 55,81\% y una precisión estratégica del 88,97\%. El modelo entrenado con 500.000 muestras mejoró estos resultados hasta el 65,64\% y el 94,01\%, respectivamente. También aumentó la precisión top-3 del 83,37\% al 92,47\% y redujo la pérdida de 1,3533 a 0,9440.

Desde el punto de vista estratégico, el incremento del conjunto de entrenamiento redujo las decisiones que empeoran el Shanten del 5,03\% al 1,59\% y las decisiones que reducen significativamente el Ukeire del 6,00\% al 4,39\%. Estos resultados, obtenidos sobre estados no utilizados durante el entrenamiento, proporcionan una evidencia más sólida de la mejora introducida por el mayor volumen de datos. No obstante, debido a la ausencia de identificadores de partida, la separación se realizó por posición en la base de datos y no pudo garantizarse la independencia entre partidas.

No obstante, la principal aportación de este trabajo no reside únicamente en el desarrollo del modelo, sino en la metodología propuesta para su evaluación. Tradicionalmente, los modelos de aprendizaje por imitación suelen evaluarse mediante métricas de clasificación, como la exactitud (\textit{accuracy}), que consideran correcta únicamente la reproducción exacta de la acción realizada por el jugador experto. En un juego como el Riichi Mahjong, donde una misma situación puede admitir múltiples decisiones estratégicamente equivalentes, este criterio resulta insuficiente y puede conducir a conclusiones erróneas sobre la calidad real del modelo.

Con el objetivo de superar esta limitación, en este trabajo se propuso una validación estratégica basada en conceptos propios del Mahjong, como el Shanten y el Ukeire. Esta metodología permite distinguir entre errores realmente perjudiciales y decisiones alternativas que mantienen o incluso mejoran la calidad estratégica de la jugada. Gracias a este enfoque fue posible observar que una parte importante de las predicciones consideradas incorrectas por las métricas tradicionales corresponden, en realidad, a decisiones estratégicamente válidas o incluso superiores a las realizadas por el jugador experto.

Sin embargo, estas métricas no deben interpretarse como una medida absoluta de la inteligencia del modelo. Una precisión estratégica elevada no implica necesariamente que el sistema haya comprendido todos los aspectos del juego, sino únicamente aquellos representados por las heurísticas utilizadas durante la evaluación. En consecuencia, estas métricas deben entenderse como una herramienta de análisis que permite describir con mayor precisión el comportamiento del modelo, y no como una medida definitiva de su nivel de juego.

\section{Líneas de Trabajo Futuro}

Aunque los resultados obtenidos son prometedores, existen numerosas líneas de trabajo que permitirían ampliar y mejorar tanto el modelo desarrollado como la metodología de evaluación propuesta.

En primer lugar, una posible línea de investigación consiste en ampliar las métricas estratégicas utilizadas durante la evaluación. En este trabajo únicamente se consideraron el Shanten y el Ukeire como indicadores de calidad estratégica. Sin embargo, el Riichi Mahjong incorpora numerosas heurísticas adicionales relacionadas con el juego ofensivo y defensivo, como el Suji, Ura-Suji, Kabe, Genbutsu o el análisis del peligro de descarte frente a jugadores en Riichi. La incorporación de este tipo de criterios permitiría construir una evaluación mucho más completa del comportamiento estratégico del modelo.

De forma similar, la propia clasificación estratégica podría extenderse para analizar otros aspectos de la toma de decisiones, permitiendo estudiar no solamente si un descarte mejora la eficiencia de la mano, sino también si representa una decisión defensivamente adecuada o consistente con estrategias de mayor nivel empleadas por jugadores expertos.

Desde el punto de vista del modelo, otra línea de trabajo consiste en explorar arquitecturas más avanzadas. Aunque en este proyecto se ha empleado un Transformer Encoder relativamente sencillo para facilitar el análisis del comportamiento del sistema, existen variantes modernas como Decision Transformer o arquitecturas específicas para aprendizaje secuencial que podrían aprovechar mejor la naturaleza temporal de las partidas de Mahjong.

Asimismo, este trabajo se ha centrado exclusivamente en la decisión de descarte, que constituye la acción más frecuente durante una partida. Como continuación natural del proyecto, sería interesante ampliar el espacio de acciones para incluir decisiones como Chi, Pon, Kan, Riichi o incluso Tsumo y Ron. Esto permitiría construir un agente capaz de participar de manera completa en una partida real de Riichi Mahjong.

Finalmente, otra línea especialmente interesante consiste en utilizar el modelo obtenido mediante aprendizaje por imitación como punto de partida para un proceso posterior de aprendizaje por refuerzo. En este contexto, las métricas estratégicas desarrolladas en este trabajo podrían resultar útiles para el diseño de funciones de recompensa más informativas que la simple victoria o derrota al final de la partida. Incorporar información relacionada con el Shanten, el Ukeire o futuras heurísticas estratégicas permitiría proporcionar una señal de aprendizaje más densa durante el entrenamiento, facilitando potencialmente la convergencia del agente y favoreciendo la adquisición de comportamientos estratégicamente sólidos.
